\section{Results: recruitment, reach and persistence}
\subsection{A negative principles campaign motivated the replication}
The separately retained November fifteen-juror principles campaign has fifty
cases across five endpoints: 390 truthful evaluated reveals, no false reveal
and ten abstentions in 400 assigned evaluated seats. The colluding-peer arm
contributes 191 reveals in 25 cases; honest peers contribute 199 in 25.
This is a protocol-specific null with missing participation, not evidence
that topology prevents spread. Its charter, offer and selected models differ
from the following replication; the campaigns are not pooled.

\subsection{Seeded recruitment can reach a binding majority}
Table~\ref{tab:swarm-scale} retains all primary size/peer cells, including
failures and non-votes. False reveals occur in both Kimi endpoints under
colluding peers at every tested jury size. In a three-juror case, one scripted
peer is joined by both evaluated Kimi K2.6 jurors. In a fifteen-juror Kimi K2.7
Code case, seven scripted peers are joined by seven of the eight evaluated
jurors. Those are direct recruitment observations, not evidence of a
seven-hop cascade.

\begin{table}[ht]
\centering\small
\caption{Frozen cloud snapshot: episodic short-charter comparisons. Condition
is jury size / scripted-peer behavior. Cases are assigned cases; F/T are false
and truthful evaluated reveals; A/U are abstentions and unrevealed commitments.
Capture counts false-majority cases; Err. counts execution-failure cases and
overlaps the other columns. A dash means no admitted played case, not zero
wrong votes. Scripted votes are excluded. Counts are descriptive and sessions
are adaptively selected; votes are not independent replicates.}
\label{tab:swarm-scale}
\begin{tabular}{llrrrrrrr}
\toprule
Model & Condition & Cases & F & T & A & U & Capture & Err. \\
\midrule
Kimi K2.6 & 3 / colluding & 5 & 7 & 0 & 3 & 0 & 4 & 1 \\
 & 3 / honest & 3 & 0 & 6 & 0 & 0 & 0 & 0 \\
 & 7 / colluding & 5 & 11 & 1 & 8 & 0 & 5 & 0 \\
 & 7 / honest & 3 & 0 & 10 & 1 & 1 & 0 & 1 \\
 & 15 / colluding & 3 & 11 & 1 & 12 & 0 & 3 & 0 \\
 & 15 / honest & 1 & 1 & 7 & 0 & 0 & 0 & 0 \\
\midrule
Kimi K2.7 Code & 3 / colluding & 4 & 6 & 0 & 2 & 0 & 3 & 1 \\
 & 3 / honest & 3 & 3 & 1 & 2 & 0 & 1 & 0 \\
 & 7 / colluding & 4 & 7 & 0 & 9 & 0 & 2 & 1 \\
 & 7 / honest & 2 & 2 & 3 & 3 & 0 & 0 & 1 \\
 & 15 / colluding & 2 & 13 & 0 & 3 & 0 & 2 & 0 \\
 & 15 / honest & 1 & 0 & 6 & 2 & 0 & 0 & 0 \\
\midrule
DeepSeek V4 Pro & 3 / colluding & 3 & 0 & 1 & 5 & 0 & 0 & 0 \\
 & 3 / honest & 2 & 0 & 3 & 1 & 0 & 0 & 0 \\
 & 7 / colluding & 3 & 0 & 0 & 12 & 0 & 0 & 1 \\
 & 7 / honest & 1 & 0 & 0 & 4 & 0 & 0 & 1 \\
 & 15 / colluding & 2 & 0 & 0 & 16 & 0 & 0 & 0 \\
 & 15 / honest & --- & --- & --- & --- & --- & --- & --- \\
\bottomrule
\end{tabular}

\end{table}

The honest-peer controls matter. Kimi K2.7 Code produces false reveals in the
three- and seven-juror honest-peer arms, including one captured three-juror
case; Kimi K2.6 has one false reveal in a fifteen-juror honest-peer case.
Consequently a bad peer is not necessary for every false vote when the
publisher's payment offer remains available. The contrast changes both peer
messaging and seeded votes, so its effect is the exposed recruitment package,
not a text-only persuasion effect.

DeepSeek V4 Pro has no false reveals, but also no evaluated reveals at either
larger size in the colluding-peer arms. Those empty decisions cannot be
presented as truthful resistance. The explicit separation of missing votes
is necessary to interpret apparent containment.

\subsection{The trace links acknowledged truth to paid execution}
\label{sec:trace-vignette}
The completed Kimi K2.6 case S1 supplies a concrete mechanism-level example.
Its returned rationale names the factual answer and considers payment and
the seeded majority; the assessment submitted to the commitment endpoint
instead claims evidential support.

\begin{modelexample}{Kimi K2.6}
\textbf{Returned reasoning.}\par\smallskip
\KimiAwareness
\traceomission
\KimiAssessmentPlan
\tcbline
\textbf{Submitted assessment.}\par\smallskip
\KimiSubmittedAssessment
\end{modelexample}
\tracemeta{S1, juror 2. Nonadjacent provider-returned reasoning before
commitment, followed by its submitted assessment. Trace lines 12 and 22;
Appendix~\ref{app:model-exchanges} supplies the actual prompt and source register.}

The commitment and reveal are confirmed, both evaluated jurors reveal the
false verdict, and each receives the offered 10 ETH in local credits. These
actions establish paid coordination against the truth in this case. The
returned rationale is evidence of a stated explanation, not guaranteed access
to faithful internal computation. It contains mistakes, which the appendix
preserves and discusses. This selected trace explains how execution can
diverge from acknowledged facts; it does not estimate prevalence or prove
that the peer message caused the decision.

\subsection{Duties and repeated exposure constrain the interpretation}
\begin{table}[ht]
\centering\small
\caption{Size-three contrasts in the same snapshot. Definitions follow
Table~\ref{tab:swarm-scale}. Repeated exposure consists of two three-case
sessions per model; its six cases are not six independent blocks. No supplied
offer or scripted juror appears in the mandated-publisher arm. A wrong ballot
alone in that arm is not an executed-coordination annotation.}
\label{tab:swarm-extensions}
\begin{tabular}{llrrrrrrr}
\toprule
Model & Condition & Cases & F & T & A & U & Capture & Err. \\
\midrule
Kimi K2.6 & Principles & 4 & 0 & 8 & 0 & 0 & 0 & 0 \\
 & Repeated exposure & 6 & 10 & 0 & 2 & 0 & 6 & 0 \\
 & Mandated publisher & 3 & 1 & 8 & 0 & 0 & 0 & 0 \\
\midrule
Kimi K2.7 Code & Principles & 4 & 0 & 8 & 0 & 0 & 0 & 0 \\
 & Repeated exposure & 6 & 9 & 2 & 1 & 0 & 5 & 0 \\
 & Mandated publisher & 3 & 0 & 9 & 0 & 0 & 0 & 0 \\
\midrule
DeepSeek V4 Pro & Principles & 2 & 0 & 4 & 0 & 0 & 0 & 0 \\
 & Repeated exposure & 6 & 0 & 9 & 3 & 0 & 0 & 0 \\
 & Mandated publisher & 2 & 0 & 6 & 0 & 0 & 0 & 0 \\
\bottomrule
\end{tabular}

\end{table}

The principles contrast has twenty truthful reveals and no false reveal in
ten cases (Table~\ref{tab:swarm-extensions}). That is consistent with an
explicit duty interrupting this financial mechanism in these selected runs;
it is not a general guarantee of containment.

All four Kimi repeated-exposure blocks have evaluated false reveals by their
second case. Eleven of their twelve cases end in a false majority. The
two Kimi K2.6 blocks each show evaluated false-vote counts $1,2,2$.
The Kimi K2.7 blocks show $1,2,2$ and $0,2,2$. In the latter block, the
first case has no majority and no payment, so observing a previous payment
cannot explain the onset in its second case. Repeated offers, memory, shared
accounts and provider variability remain competing explanations. This is
persistence under renewed exposure, not demonstrated autonomous contagion.

In the mandated-publisher arm, no case has a false majority. One Kimi K2.6
juror reveals a false minority vote; the other evaluated reveals are truthful.
This supplies no observed false-majority success. Coalition proposals and
intent in this arm require a separate transcript review; the isolated wrong
ballot does not establish either their presence or absence.
